\chapter{Metodología de la Investigación}

En este capítulo se describe el enfoque metodológico y el procedimiento seguido para diseñar, implementar, verificar y evaluar el sistema.

\section{Tipo de Investigación}

El trabajo es una investigación \textbf{aplicada y experimental}. Es aplicada porque resuelve un problema concreto ---la planificación de rutas de reparto en Cusco--- con conocimientos de algoritmia. Es experimental porque se diseñaron \textit{benchmarks} reproducibles para medir el comportamiento de cada algoritmo y contrastarlo con su análisis Big-O.

\section{Proceso Metodológico}

\subsection{Fase 1: Modelo de datos y red vial}

Se definieron las entidades en \texttt{core/modelos.py} con \texttt{dataclasses}: \texttt{Nodo} (id, nombre, latitud, longitud, sector), \texttt{Arista} (origen, destino, longitud en metros, geometría, nombres de calles y estado \texttt{activa}), \texttt{Pedido} (id, cliente, lugar de destino, coordenadas, sector, prioridad 1--3, peso, volumen y valor), \texttt{Repartidor} (vehículo, lugar de inicio y capacidades de peso y volumen) y \texttt{Ruta}. La clase \texttt{Grafo} guarda los nodos en un diccionario \texttt{id} $\to$ nodo y la adyacencia como diccionario de diccionarios.

La red vial se construyó fuera de línea con \texttt{tools/construir\_red\_vial.py}: se descargó la red de calles de Cusco de OpenStreetMap con perfil vehicular (\texttt{network\_type="drive"}, respetando sentidos de circulación), se ajustó cada uno de los 20 lugares a la intersección más cercana, y cada lugar se conectó con los 3 lugares más cercanos por calles; se añadieron 4 conexiones más para que todo lugar sea alcanzable desde cualquier otro (grafo fuertemente conexo). El peso de cada conexión es la longitud real del camino por calles. El resultado se guardó en \texttt{data/red\_vial\_cache.json} (descarga del 19 de octubre de 2026; datos © OpenStreetMap contributors, licencia ODbL), de modo que la aplicación no necesita Internet ni librerías de grafos para funcionar. Ocho lugares quedaron a más de 60 m de su coordenada original al ajustarse a la intersección más cercana (el mayor, el acceso a Sacsayhuamán, a 268 m); el Anexo B lista los ajustes.

Los datos de prueba son 30 pedidos (\texttt{data/pedidos.json}) distribuidos en los cuatro sectores, con pesos de 1 a 15 kg y prioridades 1 a 3, y 3 repartidores con bicicleta (15 kg), motocicleta (35 kg) y furgoneta (90 kg) que parten de la Plaza de Armas, Av. El Sol y Marcavalle, respectivamente.

\subsection{Fase 2: Implementación algorítmica}

Los algoritmos se implementaron en \texttt{core/}, sin interfaz gráfica, para poder probarlos de forma aislada:
\begin{itemize}
    \item \texttt{core/grafo.py}: Dijkstra con \texttt{heapq}, matriz de distancias mínimas en caché y reconstrucción de caminos.
    \item \texttt{core/estructuras.py}: Merge Sort, Quick Sort, búsqueda binaria, índice hash por ID e índice invertido por sector.
    \item \texttt{core/greedy.py}: vecino más cercano, asignación de emergencia y mochila voraz (de referencia).
    \item \texttt{core/divide\_venceras.py}: partición espacial en zonas de $K$ pedidos y en $R$ zonas, con asignación de una zona por repartidor.
    \item \texttt{core/prog\_dinamica.py}: Mochila 0/1 con tabla y reconstrucción, y Held-Karp con máscara de bits.
    \item \texttt{core/backtracking.py}: búsqueda exhaustiva de rutas simples con poda opcional por cota.
\end{itemize}
No se usaron librerías de optimización ni de grafos (\texttt{networkx}, \texttt{scipy}) en la aplicación, ni las funciones \texttt{sorted()} o \texttt{list.sort()} de Python para las ordenaciones evaluadas. Cada función documenta en su \textit{docstring} la idea del algoritmo y su complejidad temporal y espacial.

\subsection{Fase 3: Interfaz gráfica}

La GUI usa \texttt{tkinter}/\texttt{ttk} con un \texttt{Notebook} de siete pestañas: Red Vial, Pedidos, Greedy, Divide y Vencerás, Programación Dinámica, Backtracking y Comparación Big-O. El mapa usa \texttt{tkintermapview} con tiles de OpenStreetMap cuando hay conexión; si no la hay, la aplicación lo detecta al iniciar y dibuja la red en un \texttt{Canvas} vectorial propio, en el que también se pueden cerrar calles con un clic. Para cumplir el criterio de estabilidad:
\begin{itemize}
    \item Los cálculos potencialmente largos (Held-Karp, Backtracking y el \textit{benchmark} en vivo) se ejecutan en un hilo de fondo que entrega el resultado por una cola; el hilo principal la revisa cada 50 ms, de modo que ningún hilo secundario toca la interfaz (Tkinter no es seguro entre hilos) y la ventana sigue respondiendo.
    \item Toda excepción dentro de un botón se muestra en una ventana de aviso y la aplicación continúa funcionando.
    \item Las entradas del usuario se validan (peso y volumen positivos, prioridad 1--3, coordenadas en rango) y los IDs de pedidos nuevos son únicos.
    \item Las pestañas comparten un único grafo y una única lista de pedidos: al entrar a una pestaña, su mapa se redibuja si se cerró una calle en otra, y sus listas incluyen los pedidos registrados después.
    \item Los paneles laterales se desplazan verticalmente para pantallas de 1366$\times$768.
\end{itemize}

\subsection{Fase 4: Verificación y \textit{benchmarking}}

\textbf{Pruebas unitarias.} \texttt{tests/} contiene 39 pruebas \texttt{pytest} con resultados esperados conocidos (Anexo C), entre ellas: el caso clásico de mochila cuyo óptimo es 45; la comparación de la Mochila por PD contra fuerza bruta sobre los $2^{12}$ subconjuntos de 12 pedidos reales; Held-Karp sobre matrices asimétricas con óptimo calculado a mano; el Backtracking en un grafo en forma de diamante con y sin calles cerradas; y las 27 rutas de San Blas a Wanchaq.

\textbf{Prueba de la interfaz.} \texttt{tools/prueba\_gui.py} abre la aplicación y ejecuta automáticamente cada funcionalidad de cada pestaña (ordenar, buscar, registrar un pedido con datos inválidos y válidos, cada algoritmo, cerrar calles y el \textit{benchmark} en vivo), verificando que no ocurran errores ni bloqueos. Las capturas del Anexo D se generaron con esa prueba.

\textbf{\textit{Benchmarks}.} \texttt{benchmarks/benchmark.py} mide, con \texttt{time.perf\_counter()}: Merge Sort frente a Quick Sort (promedio de 20 repeticiones); Greedy frente a Held-Karp para $n\in\{3,5,8,10,12\}$; Backtracking con y sin poda en cuatro pares origen--destino; Mochila por PD frente a la voraz para $N$ de 5 a 30 pedidos y 25 kg; Divide y Vencerás frente a Greedy con 30 a 480 pedidos sintéticos (semilla fija); y la división en una zona por repartidor con los 30 pedidos reales. Los resultados se guardan en \texttt{benchmarks/resultados\_benchmark.json}, del que \texttt{informe/generar\_figuras.py} genera las figuras de este informe.

\section{Instrumentos y Entorno}

\begin{itemize}
    \item Lenguaje: Python 3.10 o superior (desarrollo y validación en Python 3.13).
    \item Dependencias de la aplicación: \texttt{tkintermapview}, \texttt{matplotlib}, \texttt{Pillow}; pruebas con \texttt{pytest}. Solo para reconstruir la red vial: \texttt{osmnx} y \texttt{networkx}.
    \item Control de versiones: Git y GitHub.
    \item Entorno de las mediciones reportadas: procesador Intel Xeon a 2.10 GHz (2 núcleos), Linux, Python 3.13. Los tiempos absolutos varían según el equipo; el orden de crecimiento y las comparaciones entre algoritmos se mantienen.
\end{itemize}
